\documentclass[12pt]{article}
\usepackage[utf8]{inputenc}
\usepackage[T1]{fontenc}
\usepackage{hyperref}
\title{Wan Language Resources for K-12 Curriculum Development}
\author{Compiled by Claude Sonnet 5 (claude-sonnet-5) for Andr\'as Kornai}
\date{}
\begin{document}
\maketitle
\section{Introduction}
Wan (ISO 639-3 wan) is a Mande language of C\^ote d'Ivoire with about 22,000 speakers (a 1993 figure), with dialects Miamu and Kemu. Web search found very little: no dictionary, primers or school materials for Wan itself. Wan should not be confused with Mwan (Mona), a separate South Mande language.
\section{Primary Educational Resources}
None found in search results.
\section{Linguistic and Reference Materials}
\begin{itemize}
\item Wan language, Wikipedia overview: \url{https://en.wikipedia.org/wiki/Wan_language}
\item Grambank entry for Wan (cites Ravenhill 1982, \emph{The Wan language}): \url{https://grambank.clld.org/languages/wann1242}
\item Numerals page for Wan\'e: \url{https://lingweb.eva.mpg.de/channumerals/Wane.htm}
\end{itemize}
\section{Religious and Community Resources}
None found in search results.
\section{Digital and Audio Resources}
None found in search results. A Lexilogos catalogue of Mande languages may point to related material: \url{https://www.lexilogos.com/mande_langues.htm}
\section{Curriculum Development Context}
Wan is not known to be taught in schools, and no curriculum materials were found. Any curriculum would have to begin with new documentation and community work.
\end{document}